% =====================================================================
%  preamble.tex - Global configuration of the Qadra user manual
%  Standalone document; it does not depend on the technical report.
%  Engine: LuaLaTeX (latexmk -lualatex main.tex)
% =====================================================================

% ---------- Language --------------------------------------------------
\usepackage[spanish,es-tabla,es-noindentfirst,es-nosectiondot]{babel}
\usepackage{csquotes}

% ---------- Fonts (same family as the technical report) ---------------
\usepackage{fontspec}
\setmainfont{Times New Roman}
\setsansfont{Times New Roman}
\setmonofont[Scale=0.88]{DejaVu Sans Mono}
\usepackage{microtype}

% ---------- Page layout -----------------------------------------------
\usepackage[letterpaper,top=2.3cm,bottom=2.3cm,left=2.5cm,right=2.3cm]{geometry}
\usepackage{setspace}
\singlespacing
\usepackage{ragged2e}
\usepackage[hyphens]{url}
\setlength{\parskip}{4pt}
\setlength{\parindent}{0pt}
% Extra stretch avoids overfull lines around long file names.
\setlength{\emergencystretch}{3em}

% ---------- Figures ---------------------------------------------------
\usepackage{graphicx}
\graphicspath{{figures/}}
\usepackage{float}
\usepackage[font=small,labelfont=bf,margin=1cm]{caption}

% ---------- Tables ----------------------------------------------------
\usepackage{booktabs}
\usepackage{longtable}
\usepackage{tabularx}
\usepackage{array}
\renewcommand{\arraystretch}{1.2}
\newcolumntype{L}[1]{>{\raggedright\arraybackslash}p{#1}}
\newcolumntype{C}[1]{>{\centering\arraybackslash}p{#1}}
\newcolumntype{Y}{>{\raggedright\arraybackslash}X}

% ---------- Lists -----------------------------------------------------
\usepackage{enumitem}
\setlist{itemsep=2pt,topsep=4pt}
\setlist[enumerate]{leftmargin=*}
\setlist[itemize]{leftmargin=*}

% ---------- Colors ----------------------------------------------------
\usepackage{xcolor}
% Main color: Qadra navy, the same as the application sidebar.
\definecolor{qadranavy}{HTML}{111344}
% Secondary color: Qadra light blue (brand-600 and brand-50 of the web design tokens).
\definecolor{qadrablue}{HTML}{2563EB}
\definecolor{qadrablue50}{HTML}{EFF6FF}
% Accent color reserved for critical warnings only (\importante).
\definecolor{ipnvino}{HTML}{6B1E3C}

% ---------- Headings --------------------------------------------------
\usepackage{titlesec}
% Parts are named A (firm staff) and B (client and verification).
\renewcommand{\thepart}{\Alph{part}}
\titleformat{\part}[display]
  {\centering\sffamily\Huge\bfseries\color{qadranavy}}
  {\Large\MakeUppercase{\partname}\ \thepart}{16pt}{}
\titleformat{\chapter}[display]
  {\sffamily\huge\bfseries\color{qadranavy}}
  {\sffamily\Large\bfseries\color{qadrablue}\MakeUppercase{\chaptertitlename}\ \thechapter}
  {12pt}{\Huge}
\titlespacing*{\chapter}{0pt}{10pt}{24pt}
\titleformat{\section}{\sffamily\Large\bfseries\color{qadranavy}}{\textcolor{qadrablue}{\thesection}}{1em}{}
\titleformat{\subsection}{\sffamily\large\bfseries}{\thesubsection}{1em}{}

% ---------- Running headers -------------------------------------------
\usepackage{fancyhdr}
\setlength{\headheight}{14pt}
\addtolength{\topmargin}{-2pt}
\pagestyle{fancy}
\fancyhf{}
\fancyhead[L]{\sffamily\small\nouppercase{\leftmark}}
\fancyhead[R]{\sffamily\small Qadra \textbullet\ Manual de usuario \quad \thepage}
\renewcommand{\headrulewidth}{0.4pt}
\renewcommand{\headrule}{{\color{qadrablue}\hrule height\headrulewidth\vskip-\headrulewidth}}
\fancypagestyle{plain}{\fancyhf{}\fancyfoot[C]{\sffamily\small\thepage}\renewcommand{\headrulewidth}{0pt}}

% ---------- Links and cross references --------------------------------
\usepackage[hidelinks,bookmarksnumbered,unicode]{hyperref}
\hypersetup{
  pdftitle={Qadra: Manual de usuario},
  pdfauthor={Eduardo Alonso Sánchez; Hatziry Vitales Herrera},
  colorlinks=true,linkcolor=qadrablue,urlcolor=qadrablue,
}
\usepackage[nameinlink,spanish]{cleveref}
\crefname{part}{parte}{partes}
\Crefname{part}{Parte}{Partes}
\crefname{chapter}{capítulo}{capítulos}
\Crefname{chapter}{Capítulo}{Capítulos}
\crefname{section}{sección}{secciones}
\Crefname{section}{Sección}{Secciones}
\crefname{subsection}{sección}{secciones}
\Crefname{subsection}{Sección}{Secciones}
\crefname{table}{tabla}{tablas}
\Crefname{table}{Tabla}{Tablas}
\crefname{figure}{figura}{figuras}
\Crefname{figure}{Figura}{Figuras}

% ---------- Manual macros ---------------------------------------------
% Placeholder for a screenshot that has not been taken yet.
% #1: path under figures/ (for example capturas/A-05-inicio-sesion.png)
% #2: what the screenshot must show and the role used to take it.
% Replace each use with \includegraphics[width=0.9\textwidth]{#1}.
\newcommand{\capturapendiente}[2]{%
  \fbox{\parbox[c][5cm][c]{0.85\textwidth}{\centering
  \textbf{CAPTURA PENDIENTE}\\[2mm]\texttt{#1}\\[1mm]\small #2}}}

% File name in monospace that may break at '-', '.', '_' and '/'.
\newcommand{\archivo}[1]{{\def\UrlBreaks{\do\/\do\-\do\.\do\_}\path{#1}}}

% Short highlighted note.
\newcommand{\nota}[1]{\par\smallskip\noindent
  \fcolorbox{qadrablue}{qadrablue50}{\parbox{\dimexpr\linewidth-2\fboxsep-2\fboxrule\relax}{%
  \small\textbf{Nota.} #1}}\par\smallskip}

% Critical warning: actions that cannot be undone or information shown once.
\newcommand{\importante}[1]{\par\smallskip\noindent
  \fcolorbox{ipnvino}{ipnvino!8}{\parbox{\dimexpr\linewidth-2\fboxsep-2\fboxrule\relax}{%
  \small\textbf{\color{ipnvino}Importante.} #1}}\par\smallskip}

% Standard block labels of each screen description.
\newcommand{\paraque}{\par\noindent\textbf{¿Para qué sirve?}\ }
\newcommand{\disponible}{\par\noindent\textbf{Disponible para:}\ }
\newcommand{\pasos}{\par\noindent\textbf{Pasos a seguir:}}
\newcommand{\resultado}{\par\noindent\textbf{Resultado esperado:}\ }
